% Tema propio de la presentación: colores, tipografía, encabezados y cajas.
% Las diapositivas están en main.tex; aquí solo va el aspecto.

\usepackage{iftex}
\ifPDFTeX
  \usepackage[utf8]{inputenc}
  \usepackage[T1]{fontenc}
\fi
% es-noshorthands: sin él, babel vuelve activos < y > y rompe las flechas de TikZ.
\usepackage[spanish,es-noshorthands,es-nolists,es-tabla,es-nodecimaldot]{babel}
\usepackage[sfdefault,lining]{FiraSans}
\usepackage[scale=0.88]{FiraMono}
\usepackage{booktabs}
\usepackage{tabularx}
\usepackage{array}
\usepackage{tikz}
\usetikzlibrary{positioning,arrows.meta,calc,babel}
\usepackage{pgfplots}
\pgfplotsset{compat=1.18}
% Números de los gráficos con la letra del texto, no con la matemática.
\pgfkeys{/pgf/number format/assume math mode=true}
% Leyenda con un solo cuadro de color por serie (el de ybar dibuja dos barras).
\pgfplotsset{leyenda cuadrada/.style={legend image code/.code={\fill[##1] (0mm,-0.9mm) rectangle (2.6mm,1.7mm);}}}
\usepackage[most]{tcolorbox}

% ---------- Paleta ----------
\definecolor{ink}{HTML}{1B2430}
\definecolor{paper}{HTML}{F6F4EF}
\definecolor{paper2}{HTML}{ECE7DC}
\definecolor{cardbg}{HTML}{FCFBF8}
\definecolor{soft}{HTML}{4A5363}
\definecolor{muted}{HTML}{5E6573}
\definecolor{rule}{HTML}{D3CCBE}
\definecolor{blue}{HTML}{2F5D9E}
\definecolor{bluesoft}{HTML}{DDE6F3}
\definecolor{bluelight}{HTML}{A9C4EA}
\definecolor{orange}{HTML}{B0501A}
\definecolor{orangesoft}{HTML}{F4E2D4}
\definecolor{orangelight}{HTML}{F0B48C}
\definecolor{night}{HTML}{243042}
\definecolor{mist}{HTML}{C9D1DC}
\definecolor{dim}{HTML}{9AA4B2}

% ---------- Beamer ----------
% Lienzo 16:9 más grande que el de Beamer (160 x 90 mm): con las mismas letras, entra más contenido.
\geometry{paperwidth=200mm, paperheight=112.5mm}
% Sin guiones al cortar palabras: en tarjetas angostas se ven mal.
\hyphenpenalty=10000
\exhyphenpenalty=10000
\setbeamercolor{background canvas}{bg=paper}
\setbeamercolor{normal text}{fg=ink}
\setbeamercolor{structure}{fg=blue}
\setbeamercolor{alerted text}{fg=orange}
\setbeamertemplate{navigation symbols}{}
\setbeamersize{text margin left=14mm, text margin right=14mm}
\setbeamertemplate{itemize item}{\color{blue}\raisebox{0.15ex}{\rule{1.1ex}{1.1ex}}}
\setbeamertemplate{itemize subitem}{\color{muted}--}
\setbeamertemplate{enumerate item}{\color{blue}\bfseries\insertenumlabel.}
\setlength{\leftmargini}{1.4em}

% Etiqueta pequeña sobre el título: la parte de la charla en la que estamos.
\newcommand{\ceja}{}
\newcommand{\parte}[1]{\renewcommand{\ceja}{#1}}

\setbeamertemplate{frametitle}{%
  \vskip 9mm
  {\color{blue}\fontsize{8}{10}\selectfont\bfseries\ceja}\par
  \vskip 1.6mm
  {\color{ink}\fontsize{20}{24}\selectfont\bfseries\insertframetitle}\par
  \vskip 1mm
}

\setbeamertemplate{footline}{%
  \hbox to \paperwidth{%
    \hspace*{14mm}%
    {\color{muted}\fontsize{6.5}{8}\selectfont Sistema Multi-Agente ReAct}%
    \hfill
    {\color{muted}\fontsize{6.5}{8}\selectfont\insertframenumber}%
    \hspace*{14mm}}%
  \vskip 5mm}

% ---------- Tipografía ----------
\newcommand{\lead}[1]{{\fontsize{13}{17}\selectfont\color{soft}#1\par}}
\newcommand{\nota}[1]{{\fontsize{9.5}{12.5}\selectfont\color{muted}#1\par}}
\newcommand{\etiqueta}[1]{{\fontsize{8}{10}\selectfont\bfseries\color{blue}#1\par}}
\newcommand{\etiquetanaranja}[1]{{\fontsize{8}{10}\selectfont\bfseries\color{orange}#1\par}}
\newcommand{\titulocaja}[1]{{\fontsize{13}{16}\selectfont\bfseries\color{ink}#1\par}}
\newcommand{\numero}[2][ink]{{\fontsize{46}{46}\selectfont\bfseries\color{#1}#2}}
\newcommand{\numerom}[2][ink]{{\fontsize{30}{30}\selectfont\bfseries\color{#1}#2}}
% Dos cifras lado a lado (multi-agente y Naive RAG), cada una sobre su etiqueta.
\newcommand{\parcifras}[2]{%
  {\renewcommand{\arraystretch}{1}%
  \begin{tabular}{@{}l@{\hspace{4mm}}l@{}}
    \numerom[blue]{#1} & \numerom[orange]{#2} \\[1mm]
    {\fontsize{8.5}{10}\selectfont\color{muted}multi-agente} & {\fontsize{8.5}{10}\selectfont\color{muted}Naive RAG}
  \end{tabular}}\par}
\newcommand{\fuerte}[1]{{\color{ink}\bfseries #1}}
\newcommand{\azul}[1]{{\color{blue}#1}}
\newcommand{\naranja}[1]{{\color{orange}#1}}

% ---------- Cajas ----------
\tcbset{
  base/.style={enhanced, boxrule=0.6pt, arc=2.2mm, left=3.2mm, right=3.2mm,
    top=2.6mm, bottom=2.6mm, before skip=0pt, after skip=0pt, halign=flush left,
    fontupper=\fontsize{10.5}{14}\selectfont\color{soft}},
}
\newtcolorbox{tarjeta}[1][]{base, colback=cardbg, colframe=rule, #1}
\newtcolorbox{tarjetaazul}[1][]{base, colback=cardbg, colframe=blue, boxrule=1pt, #1}
\newtcolorbox{tarjetanaranja}[1][]{base, colback=cardbg, colframe=orange, boxrule=1pt, #1}
\newtcolorbox{franja}[1][]{base, colback=paper2, colframe=paper2, #1}
\newtcolorbox{franjanaranja}[1][]{base, colback=orangesoft, colframe=orangesoft, #1}
\newtcolorbox{traza}[1][]{base, colback=ink, colframe=ink, left=4.5mm, right=4.5mm,
  top=3.5mm, bottom=3.5mm, fontupper=\ttfamily\fontsize{10}{15}\selectfont\color{mist}, #1}
\newtcolorbox{tarjetaoscura}[1][]{base, colback=night, colframe=night,
  fontupper=\fontsize{10.5}{14}\selectfont\color{mist}, #1}

% Líneas de la traza ReAct
\newcommand{\tpensar}[1]{#1\par}
\newcommand{\taccion}[1]{{\color{orangelight}#1}\par}
\newcommand{\tobservar}[1]{{\color{bluelight}#1}\par}
\newcommand{\tfinal}[1]{{\color{paper}\bfseries #1}\par}

% Tablas
\renewcommand{\arraystretch}{1.3}
\newcolumntype{L}{>{\raggedright\arraybackslash}X}
